\subsection*{Partie A : Phase de chauffage}
		\begin{enumerate}
			\item On a une équation différentielle de la forme \quad $y' = a y + b$,\quad où $a$ et $b$ sont des réels non nuls.

			La solution particulière constante de l'équation $(E)$ est :\quad $t \longmapsto \dfrac{-b}{a}$.

			Or \quad $\dfrac{-0,91}{-0,035} = 26$.

			La solution particulière constante de $(E)$ sur $[0 ; +\infty[$ est donc :\quad $t \longmapsto 26$.

			Par propriété, les solutions générales de $(E)$ sur $[0 ; +\infty[$ sont de la forme :\quad $t \longmapsto C\e^{-0,035t} + 26$,

			où $C$ est un réel quelconque.

			\item La fonction $T$ est donc une de ces solutions, donc on connaît sa forme générale, établie ci-avant.

			Pour déterminer la constante $C$, on va utiliser la condition initiale :\quad $T(0) = 18$.

			$\begin{aligned}
				T(0) = 18
					&\iff C\e^{-0,035\times 0} + 26 = 18\\
					&\iff C\e^{0} = 18 - 26\\
					&\iff C\times 1 = -8\\
					&\iff C = -8\\
			\end{aligned}$

			Cela confirme bien que, pour tout nombre réel $t$ appartenant à l'intervalle $[0 ;+\infty[$, on a :
			\[ T(t)=26-8 \e^{-0,035 t}.\]

			\item Selon ce modèle, le temps au bout duquel la pièce atteindra la température de 20\,\textdegree{}C est la solution de l'équation :\quad $T(t) = 20$.

			Résolvons cette équation : \quad$\begin{aligned}[t]
				T(t) = 20
					&\iff 26 - 8\e^{-0,035t} = 20\\
					&\iff -8\e^{-0,035t} = - 6\\
					&\iff \e^{-0,035t} = 0,75\\
					&\iff -0,035t = \ln(0,75)\\
					&\iff t = \dfrac{\ln(0,75)}{-0,035}
				\end{aligned}$

			$\dfrac{\ln(0,75)}{-0,035} \approx 8,22$.

			Le temps est exprimé en dizaine de minutes. 8,22 dizaines de minutes, correspondent, à la minute près à 82 minutes, soit 1 heure et 22 minutes.

			À la minute près, il faut donc 1 heure et 22 minutes pour que la température atteigne les 20\,\textdegree{}C.

			\item Pour les équations différentielles de la forme \quad $y' = ay + b$ \quad avec un coefficient $a$ strictement négatif (ce qui est le cas ici, $a = -0,035$), on sait par propriété que les courbes de fonctions vont admettre comme asymptote horizontale au voisinage de $+\infty$ la droite représentant la solution particulière constante, soit ici la droite d'équation $y = 26$.

			De plus, on sait aussi que les solutions sont strictement monotones.

			Or, ici, on a\quad $T(0) = 18 < \lim\limits_{t \to +\infty}T(t) = 26$

			donc la fonction $T$ est strictement monotone et ne peut pas être décroissante : elle est croissante (logique, en phase de chauffage, la température monte).

			$T$ est donc une fonction strictement croissante, tendant vers 26 : elle est donc majorée par 26.

			Si une panne du système de chauffage l'empêche de s'éteindre lorsque la température de la pièce atteint 20\,\textdegree{}C, la température de la pièce ne pourra donc pas dépasser 26\,\textdegree{}C, et donc encore moins 28\,\textdegree{}C.

			\emph{Remarque :} On peut aussi, plus simplement dire que la fonction exponentielle est à valeurs strictement positives et donc, pour tout réel $t$ positif :

			$\begin{aligned}
				\e^{-0,035t} > 0
					&\implies -8\e^{-0,035t} < 0\quad \text{car } -8<0\\
					&\implies 26 -8\e^{-0,035t} < 26\\
					&\implies T(t) < 26
			\end{aligned}$

			Et donc la température est majorée par 26\,\textdegree{}C, ce qui permet de conclure aussi.
		\end{enumerate}

	\subsection*{Partie B: Phase de refroidissement}
		\begin{enumerate}
			\item D'après la relation de récurrence, appliquée avec $n = 0$ :

			$u_{1}=u_{0 + 1} = 0,965u_0 + 0,35 + 0,07\e^{-0,1\times 0} = 0,965 \times 20 + 0,35 + 0,07\times 1 = 19,3 + 0,35 + 0,07 = 19,72$.

			\item Pour tout entier naturel $n$, on pose $P_n$ l'affirmation :\quad \og $u_{n}>10$\fg{}.

			\begin{itemize}[label=\textbullet~]
				\item \textbf{Initialisation :} Pour $n= 0$, on a $u_0 = 20 > 10$.

				L'affirmation $P_0$ est vraie.

				\item \textbf{Hérédité :} soit $n$ un entier naturel ($n \geqslant 0$), on suppose que pour cet entier $n$, l'affirmation $P_n$ est vraie.

				$\begin{aligned}
					P_n \text{ est vraie}
						&\implies u_n > 10\\
						&\implies 0,965 u_n > 9,65 \quad \text{car } \quad 0,965 > 0\\
						&\implies 0,965 u_n + 0,35 > 9,65 + 0,35\\
						&\implies 0,965 u_n + 0,35 > 10\\
						&\implies 0,965 u_n + 0,35 + 0,07\e^{-0,1n}> 10\quad \text{car } \quad 0,07\e^{-0,1n} > 0\\
						&\implies u_{n+1} > 10\\
						&\implies P_{n+1} \text{ est vraie}
				\end{aligned}$
				\item \textbf{Conclusion :} la propriété $P_0$ est vraie, et pour $n$ entier naturel, elle est héréditaire, donc par principe de récurrence :

				Pour tout entier naturel $n$, $P_n$ est vraie, c'est-à-dire :\quad $u_n > 10$.
			\end{itemize}

			Autrement dit la suite $(u_n)$ est strictement minorée par 10.

			\item La suite est décroissante (admis) et minorée par 10, donc d'après le théorème de convergence monotone, la suite $(u_{n})$ est convergente vers une limite supérieure ou égale à 10, notée $\ell$.

			\item \begin{enumerate}
				\item On a donc :\quad $\lim\limits_{n \to +\infty} u_n = \ell$.

				La suite de terme général $(u_{n+1})$, avec les mêmes termes, mais décalés d'un rang, aura la même limite $\ell$.

				Mais comme pour tout $n$ naturel, on a:\quad $u_{n+1} = 0,965u_n + 0,35 + 0,07\e^{-0,1n}$, on a aussi :

				\begin{itemize}[label=\textbullet~]
					\item Par limite du produit et de la somme : $\lim\limits_{n \to +\infty} 0,965u_n + 0,35 = 0,965\ell + 0,35$;
					\item $\lim\limits_{n \to +\infty} -0,1n = -\infty$\quad et donc, par composition :\quad $\lim\limits_{n \to +\infty} \e^{-0,1n} = \lim\limits_{y \to -\infty} \e^{y} = 0$;
				\end{itemize}

				Ainsi, par limite de la somme :\quad $\lim\limits_{n \to +\infty} 0,965u_n + 0,35 +0,07\e^{-0,1n} = 0,965\ell + 0,35$.

				Soit \quad $\lim\limits_{n \to +\infty} u_{n+1} = 0,965\ell + 0,35$.

				Or on a aussi \quad $\lim\limits_{n \to +\infty} u_{n+1} = \ell$.

				Par unicité de la limite, on doit donc avoir :\quad $\ell = 0,965\ell + 0,35$.

				La limite $\ell $ est donc solution de l'équation \quad $x=0,965 x+0,35$.

				\emph{Remarque :} ici, on a l'esprit du théorème du point fixe, mais on ne peut pas l'appliquer directement, car la relation de récurrence de la suite $(u_n)$ fait intervenir à la fois le terme précédent, mais aussi son indice.

				Il n'y a donc pas une fonction $f$ unique qui permettrait d'écrire $u_{n+1} = f(u_n)$, cette fonction dépendant aussi de $n$.

				\item Résolvons l'équation :\quad $\begin{aligned}[t]
					x = 0,965x + 0,35
						&\iff x - 0,965x = 0,35\\
						&\iff 0,035x = 0,35\\
						&\iff x = \dfrac{0,35}{0,035}\\
						&\iff x = 10
				\end{aligned}$

				L'équation n'a qu'une solution, 10, donc la suite tend vers 10.

				Dans le contexte de l'exercice, cela signifie qu'avec une durée sans chauffage très longue, la température de la pièce se rapproche de et se stabilise aux environs de 10\,\textdegree{}C.
			\end{enumerate}

			\item \begin{enumerate}
				\item Le programme va utiliser la variable \texttt{n} pour stocker les indices et \texttt{u} pour stocker le terme correspondant de la suite $(u_n)$.

				Tant que le terme est supérieur strictement à 18, alors la température ne nécessite pas de remettre en route le chauffage, et donc on calcule le terme et l'indice suivant.

				Cela donne :

				\begin{center}
				\begin{tblr}[T]{vline{1,3}={solid}, hline{1,Z}={solid},colspec={X[r,3mm] X[8cm,font=\ttfamily]},rowsep=0pt,column{1}={bg=gray9}}
					1.& def marche() :\\
					2.& \quad n = 0\\
					3.& \quad u = 20\\
					4.& \quad while u > 18 :\\
					5.& \quad \quad u = 0.965*u+0.35 + 0.07*exp(-0.1*n)\\
					6.& \quad \quad n = n+1\\
					7.& \quad return n\\
				\end{tblr}
				\end{center}

				\emph{Remarque :} pour que ce programme fonctionne, il faut d'abord invoquer la fonction exponentielle, fournie par la librairie \texttt{math}, avec la ligne \texttt{from math import exp}.

				\item On explore à la calculatrice les premiers termes de la suite.

				On constate que $u_8 \approx 17,87$ quand les termes précédents sont tous strictement supérieurs à 18.

				La fonction va renvoyer 8, donc il faudra huit dizaines de minutes (arrondi par excès) pour que le système de chauffage se remettra en marche.
			\end{enumerate}
		\end{enumerate}
